\Needspace{4\baselineskip}
\begin{Example}
\label{ex:ell-two-hilbert-space}
设 $\ell^2$ 为所有满足
\begin{equation}\label{eq:ell-two-sequence-square-norm}
\|x\|^2:=\sum_{n=1}^\infty|x_n|^2<\infty.
\end{equation}
的复值序列 $x_n$ 的集合. 内积定义为
\begin{equation}\label{eq:ell-two-sequence-inner-product}
(x,y)=\sum_{n=1}^\infty\overline x_ny_n.
\end{equation}
容易证明 $\ell^2$ 为内积空间. 以下证明它完备. 设
\begin{equation}\label{eq:ell-two-cauchy-sequence-components}
u^{(n)}=(u_1^{(n)},u_2^{(n)},\ldots)
\end{equation}
为 Cauchy 序列. 那么对任意 $\epsilon>0$, 存在 $N(\epsilon)$, 使
\begin{equation}\label{eq:ell-two-cauchy-sequence-tail-bound}
\|u^{(n)}-u^{(m)}\|=\sqrt{\sum_{j=1}^\infty|u_j^{(n)}-u_j^{(m)}|^2}<\epsilon
\end{equation}
对 $m,n>N(\epsilon)$ 成立. 特别地, 这意味着对每个固定的 $j$, $u_j^{(n)}$ 为 Cauchy 序列. 令
\begin{equation}\label{eq:ell-two-coordinatewise-limit}
u_j=\lim_{n\to\infty}u_j^{(n)}.
\end{equation}
由 \eqref{eq:ell-two-cauchy-sequence-tail-bound}, 可得
\begin{equation}\label{eq:ell-two-finite-partial-sum-cauchy-bound}
\sum_{j=1}^k|u_j^{(n)}-u_j^{(m)}|^2\leq\epsilon^2
\end{equation}
对每个 $n,m>N(\epsilon)$ 和每个 $k\in\mathbb{N}$ 成立. 令 $m\to\infty$, 得到
\begin{equation}\label{eq:ell-two-finite-partial-sum-limit-bound}
\sum_{j=1}^k|u_j^{(n)}-u_j|^2<\epsilon^2
\end{equation}
对 $n>N(\epsilon)$, $k\in\mathbb{N}$ 成立. 再令 $k\to\infty$, 得出 $u^{(n)}\to u$ 在 $H$ 中成立.
\end{Example}

\begin{Example}
\label{ex:lp-two-hilbert-space}
Example~\ref{ex:continuous-functions-lp-norm} 中定义的空间 $L^2(\Omega)$, 在内积
\begin{equation}\label{eq:lp-two-inner-product}
(f,g)=\int_\Omega\overline{f(\bm{x})}g(\bm{x})\,\mathrm{d}\bm{x}
\end{equation}
下为 Hilbert 空间. 这里 \eqref{eq:lp-two-inner-product} 中的积分按 Remark~\ref{rem:lp-point-values-and-measure-zero} 的意义定义.
\end{Example}

\begin{Definition}
\label{def:separable-space}
若 Hilbert 空间 (或更一般地, Banach 空间) 含有可数稠密子集, 则称其\textbf{可分}.
\end{Definition}

应用中出现的大多数空间都是可分的. 可分性对于例如通过离散化实际求解问题十分重要, 因为用这种方式只能表示空间中可数多个 (实际上, 现实世界中只有有限多个) 元素. 容易看出 $\ell^2$ 可分, 因为终止序列稠密. 空间 $L^2(\Omega)$ 也可分; 见 Problem~\ref{ex:lp-two-separable}.

\begin{Definition}
\label{def:hilbert-orthogonal-complement}
设 $H$ 为 Hilbert 空间. 若 $(x,y)=0$, 则称 $H$ 的两个元素 $x$ 和 $y$ \textbf{正交}. 对 $H$ 的任意子空间 $M$, 定义其\textbf{正交补}为
\begin{equation}\label{eq:hilbert-subspace-orthogonal-complement}
M^\perp=\{x\in H\mid(x,y)=0\ \forall y\in M\}.
\end{equation}

\end{Definition}

显然 $M^\perp$ 是闭子空间. 若 $M$ 也闭, 则 $H$ 为 $M$ 与 $M^\perp$ 的直和: $H=M\oplus M^\perp$.

\begin{Theorem}{投影定理}
\label{thm:hilbert-orthogonal-projection}
设 $H$ 为 Hilbert 空间, $M$ 为 $H$ 的闭子空间. 那么每个 $u\in H$ 都具有唯一分解 $u=v+w$, 其中 $v\in M$ 且 $w\in M^\perp$.
\end{Theorem}

\begin{Proof}
由初等几何, 预期 $v$ 是 $M$ 中离 $u$ 最近的点. 假设 $u\notin M$, 并令
\begin{equation}\label{eq:hilbert-projection-infimum-distance}
d:=\inf_{v^{\prime}\in M}\|u-v^{\prime}\|^2.
\end{equation}
那么存在序列 $v_n\in M$, 使 $d_n:=\|u-v_n\|^2$ 收敛到 $d$. 将证明 $v_n$ 为 Cauchy 序列, 并取其极限为 $v$.

设 $y$ 为 $M$ 的任意元素, $\lambda$ 为标量. 那么 $v_n+\lambda y\in M$, 因而
\begin{equation}\label{eq:projection-minimizing-sequence-perturbation}
\begin{aligned}
d&\leq\|u-(v_n+\lambda y)\|^2=\|u-v_n\|^2-\overline\lambda(y,u-v_n)\\
&\qquad-\lambda(u-v_n,y)+|\lambda|^2\|y\|^2.
\end{aligned}
\end{equation}
令 $\lambda=(y,u-v_n)/\|y\|^2$, 得到
\begin{equation}\label{eq:projection-minimizing-sequence-optimized-bound}
d\leq\|u-v_n\|^2-\dfrac{|(u-v_n,y)|^2}{\|y\|^2},
\end{equation}
它可以改写为
\begin{equation}\label{eq:projection-minimizing-sequence-pairing-bound}
|(u-v_n,y)|\leq\|y\|\sqrt{d_n-d}.
\end{equation}
接着注意
\begin{equation}\label{eq:projection-minimizing-sequence-difference-pairing}
\begin{aligned}
|(v_m-v_n,y)|&\leq|(u-v_n,y)|+|(u-v_m,y)|\\
&\leq\|y\|(\sqrt{d_n-d}+\sqrt{d_m-d}).
\end{aligned}
\end{equation}
令 $y=v_n-v_m$, 得到
\begin{equation}\label{eq:projection-minimizing-sequence-cauchy-bound}
\|v_n-v_m\|\leq\sqrt{d_n-d}+\sqrt{d_m-d},
\end{equation}
即 $v_n$ 为 Cauchy 序列. 设 $v$ 为其极限. 在 \eqref{eq:projection-minimizing-sequence-pairing-bound} 中取极限, 得到对每个 $y\in M$ 有 $(u-v,y)=0$, 即 $u-v\in M^\perp$. 这证明了所需分解的存在性.

若有两个分解 $u=v+w=v^{\prime}+w^{\prime}$, 则 $v-v^{\prime}=w^{\prime}-w$, 因而
\begin{equation}\label{eq:orthogonal-projection-decomposition-uniqueness}
\|v-v^{\prime}\|^2=(v-v^{\prime},v-v^{\prime})=(v-v^{\prime},w^{\prime}-w)=0,
\end{equation}
因为 $v-v^{\prime}\in M$ 且 $w^{\prime}-w\in M^\perp$.
\end{Proof}

\begin{Corollary}
\label{cor:hilbert-density-orthogonal-complement}
设 $H$ 为 Hilbert 空间. $H$ 的子空间 $M$ 稠密的充分必要条件为 $M^\perp=\{0\}$.
\end{Corollary}

\begin{Proof}
设 $v\in M^\perp$. 那么对每个 $u\in M$ 有 $(v,u)=0$. 若 $M$ 稠密, 则取极限可得对每个 $u\in H$ 有 $(v,u)=0$. 令 $u=v$ 得到 $v=0$.

另一方面, 由前一定理, $M$ 的闭包具有正交补. 若 $M$ 不稠密, 则这一正交补必含非零向量.
\end{Proof}
